\documentclass[11pt]{article}
\usepackage[T1]{fontenc}
\usepackage{lmodern,amsmath,amssymb,amsthm,geometry,hyperref}
\geometry{margin=1in}
\hypersetup{colorlinks=true,linkcolor=blue,citecolor=blue,urlcolor=blue,
 pdftitle={Binary consequences of relation-liveness automata lower bounds},pdfauthor={Alec Kriebel}}
\newtheorem{theorem}{Theorem}
\newtheorem{lemma}[theorem]{Lemma}
\newtheorem{proposition}[theorem]{Proposition}
\newcommand{\bits}{\{0,1\}}
\newcommand{\enc}{\operatorname{enc}}
\newcommand{\pref}{\operatorname{pref}}
\newcommand{\suff}{\operatorname{suff}}
\title{Binary consequences of relation-liveness\protect\\ automata lower bounds}
\author{Alec Kriebel\protect\\\small ORCID: \href{https://orcid.org/0009-0001-9320-500X}{0009-0001-9320-500X}}
\date{October 6, 2026}
\begin{document}
\maketitle
\begin{abstract}
We give explicit fixed-binary-alphabet consequences of OpenAI's two
relation-liveness lower bounds. Encoding each relation on $h$ points by its
$h^2$ adjacency bits gives a language recognized by a one-way nondeterministic
automaton with $N_h=(3h^3-h)/2+2$ states. Every equivalent two-way
deterministic automaton, and every two-way nondeterministic automaton for its
full binary complement, requires $2^{\Omega(h)}=2^{\Omega(N_h^{1/3})}$
states. We specify all transition rules, including an exact $sh^2$-state
pullback for an $s$-state binary target. The binary language rejects incomplete
blocks; a fooling set proves that its one-way nondeterministic state complexity
is $\Theta(h^3)$. The exponential obstructions are inherited from the two
upstream theorems; binary adjacency coding is established machinery.
\end{abstract}

\section{Scope, provenance, and conventions}
Sakoda and Sipser introduced the state cost of deterministic simulation of
two-way nondeterministic automata \cite{SakodaSipser1978}. Here we take as
inputs the two explicit relation-language lower bounds in
\cite{OAI:An-exponential-two-way-deterministic-state-lower-bound-for-one-way-liveness-September-25-2026,OAI:An-exponential-state-lower-bound-for-two-way-nondeterministic-complementation-September-25-2026}.
Both sources state their growing-alphabet scope. Our contribution is an
explicit transfer to one binary alphabet with full-word semantics, exact
state counts, and a matching cubic order for the chosen source encoding.
We do not give an independent proof of the upstream exponential obstructions.

Adjacency-bit encoding already appears in Kapoutsis's discussion of binary
one-way liveness \cite[Section~5]{Kapoutsis2011} and
\cite[Section~8]{Kapoutsis2013}, there for deterministic
targets with few reversals. That discussion does not specify the treatment
of incomplete binary blocks. Our cubic source-size claim concerns exactly
the total language defined below. Other encodings can have different source
costs. A public binary-construction record \cite{TheoremDBR816} also uses
standard coding and target pullback, with a quadratic inherited lower bound.
Neither antecedent supplies the unrestricted exponential inputs used here.
We make no claim to have invented binary coding or to have independently
resolved the base conjecture.

For comparison, Adeogun--Kapoutsis \cite[Theorem~1]{AdeogunKapoutsis2026}
proved a quadratic unrestricted deterministic lower bound for relation
liveness. Their later method limitation
\cite[Theorem~2.1]{AdeogunKapoutsisLimitation2026} bounds a particular chain
method, rather than every possible proof. Guillon--Prigioniero--Taheri
\cite[Section~1]{GuillonPrigionieroTaheri2026} describe read-only 2NFA
complementation as open; their polynomial complementation theorem uses
1-limited automata, which may rewrite cells. Our results concern read-only
machines. The upstream manuscript dates are September 25, 2026; the public
repository release was verified on October 6, 2026. These dates should not
be conflated with priority claims.

Our marked model has one initial state and a set of accepting states. The
head starts on the left endmarker of $\lhd x\rhd$; the two endmarkers are
distinct, including on empty input. A transition depends on the state and
scanned symbol, chooses a new state and a move in $\{-1,0,1\}$, and moves
the head accordingly. Outward moves at endmarkers are disabled. Partial
transitions are allowed. Acceptance is reachability of an accepting state
by a finite run, including zero transitions. Missing transitions and infinite
nonaccepting runs do not accept. A 2DFA has at most one transition per row;
a 2NFA may have several. All states are counted. Our one-way source moves
right on every body bit; its only stay is the accepting move at the right
marker. We also give an ordinary epsilon-free 1NFA, whose acceptance is
tested after consuming the whole word. The distinction matters for empty
input. Requiring a positive finite accepting run is also covered below.

\section{Language and exact source compiler}
Fix $h\ge2$, let $H=\{0,\ldots,h-1\}$, and let $\Sigma_h$ be all relations
on $H$. Products are in path order: $(p,r)\in RS$ exactly when some $q$
has $(p,q)\in R$ and $(q,r)\in S$. Write
\[
 L_h=\{R_1\cdots R_k\in\Sigma_h^*:R_1\cdots R_k\ne\varnothing\},
 \qquad R_1\cdots R_0=I_H.
\]
Set $b=h^2$. The bit at offset $t=ph+q$ of $\enc(R)$ is $1$ exactly when
$(p,q)\in R$. Extend $\enc$ by concatenation, with
$\enc(\varepsilon)=\varepsilon$, and define on \emph{every} binary word
\[
 K_h=\enc(L_h)\subseteq\bits^*.
\]
Every length-$b$ binary block codes exactly one relation. Thus $K_h$ rejects
all words whose lengths are not divisible by $b$, and otherwise tests
nonemptiness of the decoded product. It contains $\varepsilon$.

\begin{lemma}\label{lem:source}
$K_h$ has a marked one-way NFA with exactly
$N_h=(3h^3-h)/2+2$ states, and an ordinary epsilon-free 1NFA with
$N_h-1$ states.
\end{lemma}
\begin{proof}
The marked states are a nonaccepting initial state $i$, an accepting state
$a$, and
\[
 U(p,t)\quad(0\le p<h,\ 0\le t<(p+1)h),\qquad
 V(q,t)\quad(0\le q<h,\ q+1\le t<b).
\]
The offset is the next bit to be read. $U$ retains the preceding vertex
before choosing an outgoing edge; $V$ retains the chosen next vertex.
The rules at the left marker are $i\xrightarrow{\lhd,+1}U(p,0)$ for
all $p\in H$. At the right marker, precisely the states $U(p,0)$ have
the rule $U(p,0)\xrightarrow{\rhd,0}a$. State $a$ has no rules.
For a body bit $c$, give $U(p,t)$ the following rules:
\begin{itemize}
\item If $t<(p+1)h-1$, it may skip the bit and move right to $U(p,t+1)$.
\item If $c=1$ and $\lfloor t/h\rfloor=p$, put $q=t\bmod h$.
It may move right to $V(q,t+1)$ when $t<b-1$, and to $U(q,0)$ when $t=b-1$.
\end{itemize}
From $V(q,t)$ either body bit moves right to $V(q,t+1)$ for $t<b-1$,
and to $U(q,0)$ for $t=b-1$. All rules not listed are absent. Every
listed destination is in the displayed state ranges: a chosen edge at
offset $ph+q$ gives $t+1\ge q+1$. There are no left moves.

A surviving branch through a full block must choose exactly one $1$ in
the row of its preceding vertex. It then arrives at the next block as
$U(q,0)$, precisely when the block relation contains that edge. Skipping
the entire row is impossible because the last row offset has no skip rule.
Induction on complete blocks identifies surviving boundary vertices with
endpoints of paths through the decoded relations. Only complete boundaries
can accept at the right marker. This proves rejection of partial blocks
as well as correctness on valid words. On empty input, the initial right
move and the accepting stay give acceptance because $H$ is nonempty.
All accepting runs have at least two transitions, so positive-run acceptance
gives the same source language.

The two families have sizes
\[
 \sum_{p=0}^{h-1}(p+1)h=\frac{h^3+h^2}{2},\qquad
 \sum_{q=0}^{h-1}(b-q-1)=h^3-\frac{h^2+h}{2}.
\]
Adding $i,a$ proves the marked count.
For an ordinary 1NFA, retain $U,V$, replace $i,a$ by one new initial
state $i_0$, and make $i_0$ and all $U(p,0)$ final. Its bit row at $i_0$
is the union of the bit rows at $U(p,0)$ over all $p$. All other bit rows
are as above, with their right moves interpreted as ordinary consumption.
There are no transitions to $i_0$. The empty word accepts at $i_0$;
every nonempty word simulates the first bit from some starting vertex and
then the marked source's bit computation. Final states test full consumption
at a complete boundary. The count is $N_h-1$.
\end{proof}

\begin{proposition}\label{prop:cubic}
Every ordinary 1NFA recognizing $K_h$ has at least $h^3$ states. The same
bound holds for a marked right-only source in our finite-run model.
Consequently the one-way state complexity of this strict-block language
is $\Theta(h^3)$.
\end{proposition}
\begin{proof}
Let $P_p=\{(p,p)\}$ and $z=\enc(I_H)$. For $p\in H$ and $0\le t<b$, put
\[
 x_{p,t}=\enc(P_p)\pref_t(z),\qquad
 y_{p,t}=\suff_{b-t}(z)\enc(P_p).
\]
Here the prefix and suffix lengths add to $b$, including $t=0$.
For each matching pair, $x_{p,t}y_{p,t}$ encodes $P_p I_H P_p=P_p$,
so belongs to $K_h$. A crossed pair with $t\ne u$ has length
$3b+t-u$, which is not divisible by $b$, since $0<|t-u|<b$.
When $t=u$ but $p\ne r$, the crossed word encodes
$P_p I_H P_r=\varnothing$. Thus all $hb=h^3$ diagonal pairs accept and
all off-diagonal crossed pairs reject.

Choose an accepting run for each diagonal pair and its state at the cut
between $x$ and $y$. If two cuts had the same state, splicing the prefix
of one run to the suffix of the other would accept a crossed word.
The cuts must therefore occupy distinct states. This is the standard
fooling-set argument, here with its explicit witnesses.
For a marked source, an accepting run cannot accept before visiting the
right marker: otherwise the same finite prefix of its run accepts an
extension by one bit, whose length is malformed. Such a run therefore
has a cut after the prefix before reading the suffix, and the same
splicing argument applies. Stay moves, if allowed in a right-only source,
can be assigned to either side of the cut and do not affect the argument.
\end{proof}

\section{Exact pullback for all two-way targets}
\begin{lemma}\label{lem:pullback}
An $s$-state binary 2NFA $T$ has an explicit $sb$-state 2NFA $T^{\enc}$
over $\Sigma_h$ satisfying
\[
 T^{\enc}\text{ accepts }w\quad\Longleftrightarrow\quad
 T\text{ accepts }\enc(w).
\]
Determinism is preserved. This equivalence holds for both zero-step and
positive-run acceptance, with partial transitions and nonaccepting infinite
runs allowed.
\end{lemma}
\begin{proof}
If $Q$ is the state set of $T$, use $Q\times\{0,\ldots,b-1\}$.
The initial state is $(q_0,0)$; $(q,t)$ accepts exactly when $q$ accepts.
At a relation letter $R$, inspect its bit at offset $t$. For every
transition $(q,c)\mapsto(q',d)$ of $T$, where $c=\enc(R)_t$, use
\[
\begin{array}{c|c|c}
 d&\text{new offset}&\text{real move}\\\hline
 0&t&0\\
 +1, t<b-1&t+1&0\\
 +1, t=b-1&0&+1\\
 -1, t>0&t-1&0\\
 -1, t=0&b-1&-1
\end{array}
\]
with first state coordinate $q'$.
At the left marker, use $T$'s left-marker row independently of $t$:
a stay retains $t$, and a right move sets it to $0$ and moves right.
At the right marker, use its right-marker row: a stay retains $t$, and
a left move sets it to $b-1$ and moves left. Outward moves are disabled
as in $T$. No transition is supplied for a missing row.

On a relation word of length $k$, map a configuration at real body cell
$j\in\{1,\ldots,k\}$ with state $(q,t)$ to the binary configuration
with state $q$ at position $(j-1)b+t+1$. At either real marker, map all
phases to the corresponding binary marker with state $q$. The rules above
show that each transition maps to exactly one transition of $T$, and
that every transition of $T$ from the image configuration lifts from
\emph{each} corresponding pullback configuration. Marker phases need not
be unique, because marker rows ignore them. On the empty word, moves
between its adjacent markers remain correct under this mapping.
Starting from the prescribed initial states, induction on run length
proves finite-run equivalence, with the same length. Acceptance and
time-zero acceptance are preserved; hence positive-run acceptance is
preserved too. Transition-by-transition lifting also preserves infinite
runs. An infinite run supplies no new finite accepting witness on either
side. The construction has exactly $sb$ states and at most one rule per
row if $T$ does.
\end{proof}

\section{The two binary lower bounds}
We state explicitly the two external results used. For $h\ge2$, every
$s'$-state marked 2DFA for $L_h$ satisfies
\begin{equation}\label{eq:inputD}
 2^{\lfloor(h-2)/31\rfloor}\le4(s'+2)^2
\end{equation}
under either finite-run convention
\cite[Theorem~1.1]{OAI:An-exponential-two-way-deterministic-state-lower-bound-for-one-way-liveness-September-25-2026}.
Under zero-step acceptance its stronger bound uses $s'+1$.
Every $s'$-state marked 2NFA for $\Sigma_h^*\setminus L_h$ satisfies
\begin{equation}\label{eq:inputC}
 2^{\lfloor(h-2)/127\rfloor}\le2(s'+1)
\end{equation}
\cite[Theorem~1.1]{OAI:An-exponential-state-lower-bound-for-two-way-nondeterministic-complementation-September-25-2026}.
The latter source uses zero-step acceptance. Both inputs quantify over all
finite relation words and unrestricted read-only two-way targets.

\begin{theorem}\label{thm:main}
For every $h\ge2$, the explicit language $K_h\subseteq\bits^*$ of Section~2
has the source bounds of Lemma~\ref{lem:source}. If an $s$-state marked
2DFA recognizes $K_h$, then
\begin{equation}\label{eq:binaryD}
 2^{\lfloor(h-2)/31\rfloor}\le4(sh^2+2)^2,
 \qquad
 s\ge\frac{2^{\lfloor(h-2)/31\rfloor/2-1}-2}{h^2}.
\end{equation}
If an $s$-state marked 2NFA recognizes its complement in all binary words,
then
\begin{equation}\label{eq:binaryC}
 2^{\lfloor(h-2)/127\rfloor}\le2(sh^2+1),
 \qquad
 s\ge\frac{2^{\lfloor(h-2)/127\rfloor-1}-1}{h^2}.
\end{equation}
In particular, both target state costs are $2^{\Omega(N_h^{1/3})}$.
There is no polynomial state bound for either simulation or complementation
on this one fixed alphabet.
\end{theorem}
\begin{proof}
Fixed-length coding is injective on all $\Sigma_h^*$ and
$\enc(w)\in K_h$ if and only if $w\in L_h$. Apply Lemma~\ref{lem:pullback}
to a deterministic binary recognizer and insert $s'=sh^2$ into
\eqref{eq:inputD}. Taking square roots gives \eqref{eq:binaryD}.
For complementation, the full universe satisfies
\[
 \bits^*\setminus K_h
 =\{x:b\nmid |x|\}\ \cup\ \enc(\Sigma_h^*\setminus L_h),
\]
a disjoint union. The pullback sees valid codewords, and for \emph{every}
relation word $w$,
\[
 \enc(w)\in\bits^*\setminus K_h\quad\Longleftrightarrow\quad
 w\in\Sigma_h^*\setminus L_h.
\]
Thus it recognizes the entire upstream complement. Apply \eqref{eq:inputC}
with $s'=sh^2$ and rearrange. Empty input is absent from both complements;
malformed binary inputs are included in the binary complement.
Finally $h^3\le N_h\le2h^3$ for $h\ge2$. The exponential factors in
\eqref{eq:binaryD}--\eqref{eq:binaryC} absorb polynomial denominators and
constant subtractions as $h\to\infty$, proving the stated asymptotics.
The ordinary source count $N_h-1$ has the same cubic order.
\end{proof}

For a complementing target using positive-run acceptance, add a new
nonaccepting initial state with the old initial state's transition row,
whose destinations remain old states. Under zero-step acceptance this
$s+1$-state machine accepts exactly the old positive-run language:
its first transition is the required positive first step. Hence the
complementation statement remains stretched exponential, with the explicit
variant $2^{\lfloor(h-2)/127\rfloor}\le2((s+1)h^2+1)$.
The deterministic formula already covers both conventions. If a convention
requires total transition rows, an additional nonaccepting sink with legal
stay loops supplies missing rules; adding a sink changes the source bound
by one and does not weaken a lower bound for the larger partial model.
Conventions requiring acceptance at the right endmarker or halting before
accepting form subclasses of the target model; the same lower bounds apply.

The exponent here is proportional to the cube root of the binary source
size. We do not assert $2^{\Omega(n)}$ for $n$ source states, or a uniform
complexity separation such as $\mathrm L\ne\mathrm{NL}$. The transition
tables of relation-alphabet pullbacks can be large; this is a state-count
argument, not a small-table uniform algorithm claim.

\section{Verification and disclosure}
The upstream pin is \texttt{adc7f1241b42e322a6451854ab7e4b4c146bf78a}.
Independent automated audits inspected both original proof mechanisms,
the actual Lean theorem statements and machine semantics, and finite
algebraic counterexample checks. The new reduction is proved above;
its compiler and exhaustive tiny-instance simulator accompany the paper.
The simulator tests malformed lengths, empty input, legal marker behavior,
partial transitions, stays, both head directions, initial acceptance,
positive acceptance, and loops by configuration reachability. Finite tests
are falsification checks, not substitutes for uniform proofs.
A project-local attempt to rebuild the upstream Lean inputs was prevented
by insufficient disk capacity. We make no claim of reproduced formal
verification, and the binary results in this note are not formalized.
The package records exact source hashes, commands, scope and limitations.

AI tools were used extensively in research, drafting and verification.
Independent automated adversarial reviews are documented with the package.
This preprint has not undergone conventional human peer review or refereeing.
The upstream model-produced work is attributed to OpenAI as instructed by
its manuscript-specific citations. No affiliation or coauthor is asserted.

\clearpage
\begin{thebibliography}{10}
\bibitem{SakodaSipser1978} W.~J. Sakoda and M. Sipser.
Nondeterminism and the Size of Two Way Finite Automata.
STOC 1978, 275--286. \href{https://doi.org/10.1145/800133.804357}{doi:10.1145/800133.804357}.
\bibitem{OAI:An-exponential-two-way-deterministic-state-lower-bound-for-one-way-liveness-September-25-2026}
OpenAI. \emph{An exponential two-way deterministic state lower bound for
one-way liveness}. OpenAI Math Release preprint, 2026.
\href{https://github.com/openai/math/blob/adc7f1241b42e322a6451854ab7e4b4c146bf78a/preprints/An-exponential-two-way-deterministic-state-lower-bound-for-one-way-liveness-September-25-2026/main.pdf}{Pinned manuscript}, Theorem 1.1.
\bibitem{OAI:An-exponential-state-lower-bound-for-two-way-nondeterministic-complementation-September-25-2026}
OpenAI. \emph{An exponential state lower bound for two-way nondeterministic
complementation}. OpenAI Math Release preprint, 2026.
\href{https://github.com/openai/math/blob/adc7f1241b42e322a6451854ab7e4b4c146bf78a/preprints/An-exponential-state-lower-bound-for-two-way-nondeterministic-complementation-September-25-2026/paper.pdf}{Pinned manuscript}, Theorem 1.1.
\bibitem{Kapoutsis2011} C.~A. Kapoutsis. Nondeterminism Is Essential
in Small 2FAs with Few Reversals. ICALP 2011. \href{https://doi.org/10.1007/978-3-642-22012-8_15}{doi:10.1007/978-3-642-22012-8\_15}.
\bibitem{Kapoutsis2013} C. Kapoutsis. Nondeterminism Is Essential in Small
Two-Way Finite Automata with Few Reversals. \emph{Information and Computation}
222 (2013), 208--227. \href{https://doi.org/10.1016/j.ic.2012.11.001}{doi:10.1016/j.ic.2012.11.001}.
\bibitem{AdeogunKapoutsis2026} K. Adeogun and C. Kapoutsis.
A Quadratic Lower Bound for 2DFAs against One-Way Liveness.
\href{https://arxiv.org/abs/2602.24279v2}{arXiv:2602.24279v2}, July 6, 2026.
\bibitem{AdeogunKapoutsisLimitation2026} K. Adeogun and C.~A. Kapoutsis.
Unrestricted 2DFA Simulation of 1NFAs: A Quadratic Limitation to a New Lower
Bound. \href{https://arxiv.org/abs/2609.13793v2}{arXiv:2609.13793v2}, September 15, 2026.
\bibitem{GuillonPrigionieroTaheri2026} B. Guillon, L. Prigioniero and J. Taheri.
Polynomial Complementation of Nondeterministic Two-Way Finite Automata by
1-Limited Automata. STACS 2026, LIPIcs 364, 48:1--48:18.
\href{https://doi.org/10.4230/LIPIcs.STACS.2026.48}{doi:10.4230/LIPIcs.STACS.2026.48}.
\bibitem{TheoremDBR816} TheoremDB. Public record R816,
\href{https://www.theoremdb.org/records/R816/}{The one-way-liveness bound transfers to a fixed binary alphabet},
accessed October 6, 2026. Public antecedent for coding; no exponential
lower bound or peer-review status is inferred.
\end{thebibliography}
\end{document}
